\chapter{Hahn--Banach y la geometría del dual}\label{cap:v3-02}
Una funcional definida en un subespacio puede medir solo una parte del espacio. El teorema de Hahn--Banach permite extender esa medición sin empeorar su cota. Desarrollaremos primero el mecanismo de una dimensión, después la extensión global y finalmente las consecuencias para distancias, cocientes y el bidual. El ejemplo de los límites de Banach mostrará por qué una propiedad adicional debe formar parte de la definición y no suponerse automáticamente.

\section{Extender una funcional conservando una cota}\label{lec:v3-01}
\fuentelectura{01}{El teorema de Hahn--Banach}
\begin{definicion}
En un espacio vectorial real $X$, una funcional $p:X\to\R$ es sublineal si $p(x+y)\leq p(x)+p(y)$ y $p(tx)=tp(x)$ para $t\geq0$. Una funcional lineal y una norma son ejemplos. No se exige que $p$ sea no negativa ni simétrica.
\end{definicion}
\index{funcional!sublineal}
Si la homogeneidad se enuncia solo para $t>0$, tomar $x=0$, $t=2$ demuestra $p(0)=0$ y recupera el caso $t=0$. La subaditividad implica siempre $-p(-x)\leq p(x)$.

\begin{teorema}[Hahn--Banach real]\label{v3:hb-real}
Sea $M\leq X$, sea $p:X\to\R$ sublineal y sea $f:M\to\R$ lineal con $f\leq p$ en $M$. Existe una funcional lineal $F:X\to\R$ que extiende a $f$ y satisface $F\leq p$ en $X$.
\end{teorema}
\index{Hahn--Banach, teorema de}
\begin{proof}
\emph{Paso de una dimensión.} Si $x_0\notin M$, la representación $m+tx_0$ en $M+\R x_0$ es única: dos representaciones darían un múltiplo no nulo de $x_0$ en $M$. Toda extensión lineal tiene entonces la forma $f_c(m+tx_0)=f(m)+tc$. Para que $f_c\leq p$, los casos $t>0$ y $t<0$ requieren, respectivamente,
\[
 c\leq p(z+x_0)-f(z)\quad(z\in M),\qquad
 c\geq f(y)-p(y-x_0)\quad(y\in M).
\]
En efecto, divida por $t$ en el primer caso y escriba $t=-s$, $s>0$, en el segundo. Para $t=0$ ya tenemos la hipótesis original. El intervalo de valores admisibles es no vacío porque
\[
 f(y)+f(z)=f(y+z)\leq p(y+z)\leq p(y-x_0)+p(z+x_0).
\]
Así cada cota inferior es menor o igual que cada cota superior. El supremo de las inferiores y el ínfimo de las superiores son reales: tomar $y=0$ y $z=0$ da cotas finitas en ambos sentidos. Elija $c$ entre ellos.

\emph{Paso global.} Considere los pares $(N,g)$ donde $M\subseteq N\leq X$, $g$ extiende a $f$ y $g\leq p$ en $N$. Ordénelos por extensión. El conjunto es no vacío porque contiene $(M,f)$. En una cadena, las funciones coinciden donde se superponen sus dominios. La unión es, por tanto, una función bien definida. Su dominio es un subespacio: dos vectores de la unión pertenecen juntos al dominio mayor de dos miembros comparables, y sus combinaciones lineales también. Esa misma observación verifica linealidad; cada valor sigue dominado por $p$ y la función sigue extendiendo a $f$. La unión es una cota superior de la cadena. La cadena vacía tiene, por ejemplo, la cota $(M,f)$.

El lema de Zorn proporciona un par maximal $(N,g)$. Si $N\ne X$, el paso de una dimensión construye una extensión propia dominada por $p$, contradicción. Por tanto $N=X$ y $F=g$.
\end{proof}

\begin{teorema}[Extensión real y compleja sin aumentar la norma]\label{v3:hb-norma}
Sea $M$ un subespacio de un espacio normado real o complejo $X$. Cada $f\in M^*$ admite una extensión $F\in X^*$ con $\norm F=\norm f$. No se exige que $M$ sea cerrado ni que $X$ sea completo.
\end{teorema}
\begin{proof}
En el caso real, aplique el teorema anterior a $p(x)=\norm f\norm x$. La cota $F(x)\leq\norm f\norm x$, aplicada también a $-x$, da $|F(x)|\leq\norm f\norm x$. Así $\norm F\leq\norm f$, mientras que restringir a $M$ da la desigualdad contraria.

En el caso complejo, considere los espacios reales subyacentes y extienda $\operatorname{Re}f$ a una funcional real $G$ con $G(x)\leq\norm f\norm x$. Defina
\[
 F(x)=G(x)-iG(ix).
\]
Es aditiva y real homogénea. Además $F(ix)=G(ix)+iG(x)=iF(x)$, de modo que es compleja lineal. Para $m\in M$, $G(m)=\operatorname{Re}f(m)$ y $G(im)=-\operatorname{Im}f(m)$; por tanto $F(m)=f(m)$. Si $F(x)\ne0$, elija $|\lambda|=1$ con $\lambda F(x)=|F(x)|$. Entonces
\[
 |F(x)|=\operatorname{Re}F(\lambda x)=G(\lambda x)\leq\norm f\norm x.
\]
El caso $F(x)=0$ es inmediato. Esto prueba $\norm F\leq\norm f$; la restricción vuelve a dar igualdad. La lectura atribuye esta versión compleja a Bohnenblust, Sobczyk y Suhomlinov; aquí se distingue el argumento real del paso adicional complejo.
\end{proof}
El operador de restricción $X^*\to M^*$ es lineal, contractivo y sobreyectivo. La cota contractiva se debe a que el supremo se toma sobre una bola menor. Puede tener núcleo no trivial: en $\R^2$, la segunda coordenada se anula en el eje horizontal.

\section{Norma, distancia y bidual}\label{lec:v3-02}
\fuentelectura{02}{Primeras aplicaciones del teorema de Hahn--Banach}
\begin{proposicion}[Funcionales normantes]\label{v3:normantes}
Para cada $x\ne0$ existe $\varphi\in X^*$ con $\norm\varphi=1$ y $\varphi(x)=\norm x$. En consecuencia,
\[
 \norm x=\sup_{\norm\varphi\leq1}|\varphi(x)|.
\]
El dual separa puntos y la aplicación canónica $J_X:X\to X^{**}$, $J_Xx(\varphi)=\varphi(x)$, es una isometría lineal.
\end{proposicion}
\begin{proof}
En $\F x$ defina $f(\alpha x)=\alpha\norm x$. Tiene norma uno; extiéndala por Hahn--Banach. Toda funcional de norma a lo sumo uno satisface $|\varphi(x)|\leq\norm x$, y la construida alcanza esa cota. Para $x=0$ la fórmula es directa. Si $u\ne v$, una funcional normante de $u-v$ satisface $\varphi(u)-\varphi(v)=\norm{u-v}\ne0$. La evaluación $J_Xx$ es lineal sobre $X^*$ y acotada por $\norm x$; la fórmula del supremo demuestra $\norm{J_Xx}=\norm x$. Su linealidad respecto de $x$ se verifica evaluando en cualquier $\varphi$.
\end{proof}
\index{bidual}\index{funcional!normante}
No se afirma que toda funcional alcance su norma. Para cada vector se construye una funcional adecuada; intercambiar esas dos elecciones produciría una afirmación distinta y generalmente falsa.

\begin{proposicion}[Medir la distancia a un subespacio]\label{v3:distancia-dual}
Si $M\leq X$ y $d=\dist(x,M)>0$, existe $\varphi\in X^*$ con $\norm\varphi=1$, $\varphi|_M=0$ y $\varphi(x)=d$.
\end{proposicion}
\begin{proof}
Como $x\notin M$, cada vector de $M+\F x$ tiene una única representación $m+\alpha x$. Defina $f(m+\alpha x)=\alpha d$. Si $\alpha\ne0$,
\[
 |f(m+\alpha x)|=|\alpha|d\leq|\alpha|\norm{x+m/\alpha}=\norm{m+\alpha x};
\]
si $\alpha=0$, la desigualdad también vale. Así $\norm f\leq1$. Elija $m_n\in M$ con $\norm{x-m_n}\to d$. Entonces $d=|f(x-m_n)|\leq\norm f\norm{x-m_n}$ y, al pasar al límite y dividir por $d>0$, $\norm f\geq1$. Extender $f$ conserva su norma y sus valores.
\end{proof}

\begin{definicion}
Para $S\subseteq X$ y $A\subseteq X^*$ definimos el anulador y preanulador
\[
 S^\perp=\{\varphi\in X^*: \varphi(s)=0\ \forall s\in S\},\qquad
 {}^\perp A=\{x\in X:\varphi(x)=0\ \forall\varphi\in A\}.
\]
Aquí $\perp$ no denota necesariamente ortogonalidad de Hilbert; indica anulación bajo el apareamiento entre un espacio y su dual.
\end{definicion}
\index{anulador}\index{preanulador}
Son subespacios cerrados: cada uno es una intersección de núcleos de funciones lineales continuas, en el segundo caso evaluaciones. La continuidad también da $M^\perp=(\cl M)^\perp$. Para todo subespacio $M$,
\begin{equation}\label{v3:doble-anulador}
 {}^\perp(M^\perp)=\cl M.
\end{equation}
Una inclusión es inmediata. Si $x\notin\cl M$, la proposición de distancia proporciona una funcional que se anula en $M$ y no en $x$, lo que excluye a $x$ del preanulador. Esta es también la demostración completa de la proposición8 que la lectura deja como ejercicio.

\begin{proposicion}[Sistema biortogonal finito]\label{v3:biortogonal}
Dados $x_1,\ldots,x_m$ linealmente independientes, existen $\varphi_1,\ldots,\varphi_m\in X^*$ con $\varphi_i(x_j)=\delta_{ij}$. Para $x$ en su espacio generado, $x=\sum_j\varphi_j(x)x_j$.
\end{proposicion}
\begin{proof}
En $E=\spanop(x_1,\ldots,x_m)$, las funcionales coordenada son continuas por equivalencia de normas en dimensión finita. Extiéndalas por Hahn--Banach. La fórmula se verifica escribiendo $x=\sum_ja_jx_j$ y aplicando cada coordenada. Alternativamente, se puede aplicar la proposición de distancia a $x_i$ y al espacio generado por los demás vectores, y dividir la funcional resultante por la distancia positiva; ese factor es necesario para obtener valor uno.
\end{proof}

\section{Ejercicios propuestos: extensión algebraica y sublinealidad}
\begin{ejercicio}[Lectura1, ejercicio1]\label{ej:v3-l01-1}
Si $p,q$ son sublineales reales y $\alpha>0$, pruebe que $p+q$ y $\alpha p$ son sublineales.
\end{ejercicio}

\begin{ejercicio}[Lectura1, ejercicio2]\label{ej:v3-l01-2}
Sea $(p_n)$ creciente y sublineal, con límite real $p(x)$ para cada $x$. Pruebe que $p$ es sublineal.
\end{ejercicio}

\begin{ejercicio}[Lectura1, ejercicio3]\label{ej:v3-l01-3}
Dados espacios vectoriales reales $X,Y$, un subespacio $M$ y una aplicación lineal $T:M\to Y$, construya una extensión lineal a $X$.
\end{ejercicio}

\begin{ejercicio}[Lectura1, ejercicio4]\label{ej:v3-l01-4}
Construya una solución de $f(x+y)=f(x)+f(y)$ en $\R$ que no tenga la forma $f(x)=cx$.
\end{ejercicio}


\section{Límites de Banach y convergencia casi uniforme de promedios}
\index{Banach, limite de@Banach, límite de}\index{casi convergencia}
La funcional límite en $c\subseteq\ell^\infty$ tiene norma uno: $|\lim x_n|\leq\norm x_\infty$, y la sucesión constante uno alcanza la cota. Hahn--Banach produce extensiones de norma uno a $\ell^\infty$. Sin embargo, \emph{esas extensiones no son automáticamente invariantes por desplazamientos}. La denominación de la lectura necesita reforzarse antes de resolver su ejercicio5.

\begin{definicion}[Convención corregida]
Un límite de Banach real es una funcional lineal $B:\ell^\infty(\R)\to\R$ positiva, con $B(\mathbf1)=1$ e invariante por el desplazamiento izquierdo $S(x_1,x_2,\ldots)=(x_2,x_3,\ldots)$. Su extensión compleja se define por $B(x+iy)=B(x)+iB(y)$ para sucesiones reales $x,y$.
\end{definicion}
La positividad implica $|B(x)|\leq\norm x_\infty$ para $x$ real. Además $B(e_1)=B(Se_1)=0$ y, por desplazamiento, $B(e_n)=0$ para todo $n$; por continuidad, $B$ se anula en $c_0$. Por tanto extiende el límite usual. También $B(Rx)=B(x)$ para el desplazamiento derecho $R$, pues $SR=I$. En el caso complejo, elija una fase $\lambda$ que haga real no negativo $\lambda B(z)$; entonces $|B(z)|=B(\operatorname{Re}(\lambda z))\leq\norm z_\infty$. La extensión tiene norma uno.

\begin{lema}\label{v3:banach-limites}
Para $x\in\ell^\infty(\R)$ defina
\[
 p(x)=\lim_{k\to\infty}\sup_{n\geq1}\frac1k\sum_{j=0}^{k-1}x_{n+j},\qquad q(x)=-p(-x).
\]
Estos límites existen. La función $p$ es sublineal y los valores posibles de $B(x)$, al variar los límites de Banach, forman exactamente $[q(x),p(x)]$.
\end{lema}
\begin{proof}
Ponga $a_k=\sup_n\sum_{j=0}^{k-1}x_{n+j}$. Separar un bloque en dos da $a_{k+l}\leq a_k+a_l$. Para un $m$ fijo escriba $k=tm+r$, $0\leq r<m$. Entonces $a_k\leq ta_m+r\norm x_\infty$. Al dividir por $k$ y tomar límite superior obtenemos $\limsup a_k/k\leq a_m/m$. Como cada $a_k/k$ es al menos $\inf_m a_m/m$, existe el límite y es ese ínfimo. Es finito, entre $-\norm x_\infty$ y $\norm x_\infty$.

Cada promedio es lineal; tomar supremos y después límites da $p(x+y)\leq p(x)+p(y)$ y homogeneidad positiva. También $p(x+c\mathbf1)=p(x)+c$. Los promedios de $Sx-x$ son $(x_{n+k}-x_n)/k$, de módulo a lo sumo $2\norm x_\infty/k$; por ello $p(Sx-x)=p(x-Sx)=0$.

Hahn--Banach aplicado a $c\mathbf1\mapsto c$, dominado por $p$, da una funcional $B\leq p$. Aplicando la dominación a $-x$ se obtiene $q(x)\leq B(x)$. Si $x\geq0$, $q(x)\geq0$ y $B$ es positiva. Las dos cotas sobre $Sx-x$ dan invariancia. Así existen límites de Banach.

Para alcanzar un valor $c\in[q(x),p(x)]$, defina primero $f(tx+s\mathbf1)=tc+s$. Si $x$ no es constante, la representación es única; si es constante, el intervalo contiene solo esa constante y la definición sigue siendo consistente. Para $t\geq0$, $tc+s\leq tp(x)+s=p(tx+s\mathbf1)$. Para $t<0$, $c\geq-p(-x)$ da la misma dominación. Hahn--Banach extiende $f$ dominada por $p$, y el argumento precedente produce un límite de Banach con $B(x)=c$.

Finalmente, todo límite de Banach satisface $B(x)=B(k^{-1}\sum_{j=0}^{k-1}S^jx)$. La positividad lo coloca entre el ínfimo y el supremo de las entradas de ese promedio. Al pasar $k\to\infty$, $q(x)\leq B(x)\leq p(x)$.
\end{proof}

\begin{ejercicio}[Lectura1, ejercicio5, incisos(i)--(ii)]\label{ej:v3-l01-5}
Una sucesión es casi convergente si todos sus límites de Banach tienen el mismo valor. (i) Estudie la sucesión $(1,0,1,0,\ldots)$. (ii) Pruebe que $x$ es casi convergente a $\alpha$ si y solo si sus promedios $k^{-1}\sum_{j=0}^{k-1}x_{n+j}$ convergen a $\alpha$ uniformemente respecto de $n$. Distinga la definición corregida de la extensión arbitraria utilizada inicialmente en la lectura.
\end{ejercicio}


\section{Ejercicios propuestos: anulación y cocientes}
\begin{ejercicio}[Lectura2, ejercicio1]\label{ej:v3-l02-1}
Si $X^*$ es separable, demuestre que $X$ es separable.
\end{ejercicio}

\begin{ejercicio}[Lectura2, ejercicio2, incisos(i)--(iv)]\label{ej:v3-l02-2}
Pruebe las identidades y contenciones: (i) $\{0\}^\perp=X^*$ y ${}^\perp X^*=\{0\}$; (ii) $M_1\subseteq M_2$ implica $M_1^\perp\supseteq M_2^\perp$; (iii) $(S\cup T)^\perp=S^\perp\cap T^\perp$ y $S^\perp+T^\perp\subseteq(S\cap T)^\perp$; (iv) $S\subseteq{}^\perp A$ equivale a $\varphi(s)=0$ para todo $s\in S$, $\varphi\in A$, y a $A\subseteq S^\perp$.
\end{ejercicio}

\begin{ejercicio}[Lectura2, ejercicio3: examen del enunciado]\label{ej:v3-l02-3}
El original pide, para $M_1\cap M_2=\{0\}$ finito dimensionales, funcionales $\varphi_i$ que se anulen en $M_i$ y sean no nulas en todo vector del otro subespacio. Determine cuándo es posible y formule una versión válida para dimensiones arbitrarias.
\end{ejercicio}

\begin{ejercicio}[Lectura2, ejercicio4]\label{ej:v3-l02-4}
Para $x\in X$, pruebe que $\{\varphi\in X^*: \varphi(x)=\norm x\}$ es no vacío y convexo.
\end{ejercicio}

\begin{ejercicio}[Lectura2, ejercicio5]\label{ej:v3-l02-5}
Si $M$ es un subespacio cerrado de $X$, pruebe que
\[
 \Phi:M^\perp\to(X/M)^*,\quad\Phi(\varphi)(x+M)=\varphi(x),\qquad
 \Psi:X^*/M^\perp\to M^*,\quad\Psi(\varphi+M^\perp)=\varphi|_M
\]
son isomorfismos. Precise sus normas.
\end{ejercicio}


\section{Del álgebra a la geometría}
Hahn--Banach convierte un cálculo en un subespacio pequeño en una funcional sobre todo el espacio. De ahí proceden la fórmula dual de la norma, la recuperación de subespacios cerrados mediante anuladores y las identificaciones de duales de cocientes. La extensión no añade propiedades que no se hayan impuesto: el caso de los límites de Banach distingue con precisión extensión, positividad e invariancia. En el capítulo siguiente usaremos funciones sublineales para producir hiperplanos que separan conjuntos convexos.
